%% 트레이닝복 판: figures/cogsci_arena_round.tex 복사본에서 GIVE·ASK·EXIT 하는 로봇만 포즈 스프라이트로 바꿈.
%% 공존 아레나 v33의 한 라운드 (2026-10-20 구현 기준), 오징어 게임 장면처럼 다시 그린 판(2026-10-21).
%% 라운드 시작 네 자리이므로 D의 EXIT 뒤에도 이번 벌금은 24,000/4 = 6,000.
%% 끝 할당은 마지막까지 참여한 N_end 자리에 총 20,000,000을 균등 분할한다.
%% 규칙은 캡션과 본문이 설명하므로 그림은 흐름만 보인다. 허용한 글: 단계 머리, A-D,
%% GIVE / ASK / EXIT / BELIEF 표, -6,000과 0, 고리 위 "다음 라운드"와 끝 할당.
%% 이야기: 1 높은 발판의 진행요원이 총으로 B를 겨눠 solver로 고른다. 시작 막대 4/3/2/1칸 = 60/40/20/10%.
%%         2 B만 말풍선(토큰 ?)으로 ASK, A는 말 없이 GIVE. 준 쪽(A) 막대의 넘긴 칸은 빗금,
%%           받은 쪽(B) 막대에는 준 쪽의 색(파랑)으로 칸을 채운다. 각 자리 밑 EXIT 표지(여기서는 나갈 수 있음).
%%           D는 EXIT 표지가 달린 문으로 걸어 나간다.
%%         3 B가 카드를 들고 풀며 BELIEF 막대를 낸다(토큰이 줄어 1칸).
%%         4 위쪽을 둘로 나눔: 정답(체크)이면 B에게 동전이 돌아옴 / 오답(X)이면 진행요원이 -6,000 팻말로
%%           벌금을 알림. A가 곧 빌 C에게 GIVE(A는 빗금, C는 파랑 칸). 각 자리 밑 EXIT 표지에 X.
%%         5 벌금을 한꺼번에 뗌(빨간 칸): A와 C는 남고 SOLVE에 토큰을 쓴 B는 0으로 꺼짐.
%% 논문판(2026-11-05): 단계 2~5를 xshift scope로 밀어 단계 사이를 넓힘(아래 \sgdB~\sgdE).
%% 3차: 에이전트 알파벳 라벨을 지움(색으로 구분), ASK를 말풍선 왼쪽 같은 높이로, 인물·막대·소품을 키움(\sgS·\sgY·\sgB).
%% 조각은 cogsci2027/figpics.tex. 이전 판: arena_round_v1.tex(부록에 실음).
\begin{tikzpicture}[x=1cm,y=1cm,
  lab/.style={sgfont=6.6,inner sep=0pt,text=black},
  tag/.style={sgfont=6.8,inner sep=1.8pt,rounded corners=1.5pt,
    font=\sffamily\bfseries\fontsize{6.8}{7}\selectfont},
  fine/.style={tag,draw=sgastra,fill=sgastra!10,text=black},
  spot/.style={fill=sggold!25},
  laser/.style={sgastra,line width=.55pt,dash pattern=on 2pt off 1pt}]

%% == 단계 사이 간격 (2026-11-05 논문판) ==
%% 단계 2~5를 오른쪽으로 \sgdB~\sgdE만큼 민다. 합 3.3은 본문에서 그림 높이를 그대로 두고 폭을 채우는 양.
%% 2와 3 사이(문과 조명 사이 0.24)가 가장 좁았으므로 가장 많이 벌린다. 꺾쇠는 두 단계 내용의 가운데.
\pgfmathsetmacro\sgdB{.75}
\pgfmathsetmacro\sgdC{2.0}
\pgfmathsetmacro\sgdD{2.65}
\pgfmathsetmacro\sgdE{3.3}
\pgfmathsetmacro\sgW{17.8+\sgdE}
%% == 크기 (2026-11-05 논문판 3차): 알파벳 라벨을 지운 자리만큼 인물을 키우고 올린다 ==
%% \sgS 인물 배율(원래 1.25), \sgY 인물 기준선(원래 2.0), \sgB 토큰 막대 배율(원래 1.45).
\def\sgS{1.5}
\def\sgY{2.12}
\def\sgB{1.6}
%% == 단계 머리 ==
\sgstep{(1.6,4.5)}{sgblue}{1}{푸는 에이전트}
\sgstep{(5.4+\sgdB,4.5)}{sgpurple}{2}{풀기 전 대화}
\sgstep{(9.05+\sgdC,4.5)}{sgorange}{3}{풀기}
\sgstep{(12.7+\sgdD,4.5)}{sggreen}{4}{결과 뒤 대화}
\sgstep{(16.35+\sgdE,4.5)}{sgastra}{5}{벌금}
\draw[black!20,line width=.5pt] (0,4.22)--(\sgW,4.22);
%% 단계 내용의 좌우 끝: 1 오른쪽 3.2 / 2 왼쪽 3.6, 오른쪽 7.48 / 3 왼쪽 7.61, 오른쪽 10.2 / 4 왼쪽 10.7, 오른쪽 14.45 / 5 왼쪽 15.0
\sgchev{({(3.2+3.6+\sgdB)/2},2.2)}
\sgchev{({(7.48+\sgdB+7.61+\sgdC)/2},2.2)}
\sgchev{({(10.2+\sgdC+10.7+\sgdD)/2},2.2)}
\sgchev{({(14.45+\sgdD+15.0+\sgdE)/2},2.2)}

%% == 1. solver: 발판 위 진행요원이 B를 겨눔, 시작 막대 4/3/2/1 ==
\fill[spot] (1.1,4.0)--(1.3,4.0)--(1.7,1.12)--(.7,1.12)--cycle;
\fill[black!35] (2.2,2.86) rectangle (3.2,2.93);
\draw[black!35,line width=.6pt] (2.28,2.86)--(2.28,2.68) (3.12,2.86)--(3.12,2.68);
\sgguardgun{(2.85,2.93)}{38}
\draw[laser] (sgmuzzle)--(1.37,2.19);
\draw[sgastra,line width=.5pt] (1.37,2.19) circle (.06);
\fill[sgastra] (1.37,2.19) circle (.025);
\sgrobot[\sgS]{(.4,\sgY)}{sgblue}{}
\sgrobot[\sgS]{(1.2,\sgY)}{sgorange}{}
\sgrobot[\sgS]{(2.0,\sgY)}{sggreen}{}
\sgrobot[\sgS]{(2.8,\sgY)}{sgpurple}{}
\sgbar[\sgB]{(.4,1.3)}{sgblue}{4}{4}
\sgbar[\sgB]{(1.2,1.3)}{sgorange}{3}{4}
\sgbar[\sgB]{(2.0,1.3)}{sggreen}{2}{4}
\sgbar[\sgB]{(2.8,1.3)}{sgpurple}{1}{4}
\foreach \x/\p in {.4/60,1.2/40,2.0/20,2.8/10}
  \node[lab] at (\x,1.03) {\p k};

%% == 2. TALK: B만 ASK(말풍선), A는 말 없이 GIVE / 각자 밑 EXIT 표지 / D는 문으로 나감 ==
\begin{scope}[xshift=\sgdB cm]
\sgbubble[inner sep=2.6pt]{(5.66,3.2)}{(5.42,2.82)}{%
  \tikz[baseline=-2.5pt]{\sgtoken[1.05]{(0,0)}}\,{\sffamily\bfseries\fontsize{9}{9}\selectfont ?}}
\node[tag,draw=sgorange,fill=sgorange!12,text=black] at (4.92,3.2) {ASK};  % 말풍선과 0.12 띄움
\sgpose[\sgS]{(4.0,\sgY)}{give_blue}{}
\sgpose[\sgS]{(5.35,\sgY)}{ask_orange}{}
\sgbar[\sgB]{(4.0,1.3)}{sgblue}{3}{4}
\sgused[\sgB]{(4.0,1.3)}{sgblue}{4}{4}
\sgbar[\sgB]{(5.35,1.3)}{sgorange}{4}{4}
\sggotfrom[\sgB]{(5.35,1.3)}{sgblue}{1}{4}
\draw[sgflow,sgblue!75] (4.36,1.3)--(4.99,1.3);
\sgcoins[1.05]{(4.67,1.37)}{2}
\node[tag,draw=sgblue,fill=sgblue!10,text=black] at (4.67,1.84) {GIVE};
\sgexitsign[1.0]{(4.0,.86)}
\sgnoexitsign[1.0]{(5.35,.86)}
\sgexitdoor[1.35]{(6.8,1.17)}
\sgexitsign[1.1]{(6.8,2.78)}
\sgpose[1.12]{(6.88,1.72)}{walk_purple}{}
\draw[sgpurple!60,line width=.6pt,line cap=round] (6.3,1.92)--(6.47,1.92) (6.27,1.77)--(6.47,1.77) (6.3,1.62)--(6.47,1.62);
\sgbar[1.3]{(6.8,.98)}{sgpurple}{1}{4}
\node[lab,text=black] at (6.8,.76) {+20,000};
\end{scope}

%% == 3. SOLVE: 조명 아래 B가 카드를 들고 BELIEF를 냄 ==
\begin{scope}[xshift=\sgdC cm]
\fill[spot] (7.98,4.0)--(8.18,4.0)--(8.55,1.12)--(7.61,1.12)--cycle;
\sgrobot[\sgS]{(8.08,\sgY)}{sgorange}{}
\sgbar[\sgB]{(8.08,1.3)}{sgorange}{1}{4}
\begin{scope}[shift={(9.1,1.8)},scale=1.2]
  \filldraw[draw=black!30,fill=white,rounded corners=2pt] (-.4,-.3) rectangle (.48,.38);
  \filldraw[draw=black!50,fill=white,rounded corners=2pt,line width=.6pt] (-.44,-.34) rectangle (.44,.34);
  \fill[sgorange] (-.32,-.14)--(-.18,.14)--(-.04,-.14)--cycle;
  \fill[sggreen] (.11,.02) circle (.11);
  \fill[sgblue] (.25,-.14) rectangle (.36,-.03);
  \fill[sgpurple] (.25,.06) rectangle (.36,.17);
\end{scope}
\begin{scope}[shift={(9.475,3.32)},scale=1.15,shift={(-9.475,-3.32)}]
\fill[black!25] (8.52,2.8) circle (.045) (8.68,2.96) circle (.06);
\filldraw[draw=black!45,fill=white,rounded corners=4pt,line width=.5pt] (8.85,2.86) rectangle (10.1,3.78);
\begin{scope}[shift={(9.07,3.0)}]
  \foreach \i/\h in {0/.14,1/.26,2/.62,3/.18}
    \filldraw[draw=sgorange!80!black,fill=sgorange!60,line width=.35pt] ({\i*.22},0) rectangle ++(.15,\h);
  \draw[black!50,line width=.45pt] (-.05,0)--(.86,0);
\end{scope}
\end{scope}
\node[tag,draw=sgorange,fill=sgorange!12,text=black] at (9.48,2.62) {답};
\end{scope}

%% == 4. 결과 + TALK: 맞히면 B에게 환급 / 틀리면 진행요원이 -6,000을 알림, A가 C에게 GIVE ==
\begin{scope}[xshift=\sgdD cm]
\sgcheck[1.1]{(10.98,3.66)}
\node[lab,text=black] at (10.98,3.27) {정답};
\sgcoins[1.05]{(11.42,3.54)}{3}
\draw[sgflow,sggold!80!black] (11.62,3.66)--(11.95,3.66);
\sgrobot[.62]{(12.2,3.6)}{sgorange}{}
\draw[black!20,line width=.5pt] (12.58,3.98)--(12.58,3.12);
\sgcross[1.1]{(12.95,3.66)}
\node[lab,text=black] at (12.95,3.27) {오답};
\sgguardsign[.8]{(14.0,2.5)}{1.5}{$-$6,000}  % 팻말 위끝 4.04: 머리 구분선(4.22)과 띄움
\sgpose[\sgS]{(11.05,\sgY)}{give_blue}{}
\sgrobot[\sgS]{(12.42,\sgY)}{sggreen}{}
\sgrobot[\sgS]{(13.45,\sgY)}{sgorange}{}
\sgbar[\sgB]{(11.05,1.3)}{sgblue}{2}{4}
\sgused[\sgB]{(11.05,1.3)}{sgblue}{3}{4}
\sgbar[\sgB]{(12.42,1.3)}{sggreen}{3}{4}
\sggotfrom[\sgB]{(12.42,1.3)}{sgblue}{1}{4}
\sgbar[\sgB]{(13.45,1.3)}{sgorange}{1}{4}
\draw[sgflow,sgblue!75] (11.39,1.3)--(12.08,1.3);
\sgcoins[1.05]{(11.73,1.37)}{2}
\node[tag,draw=sgblue,fill=sgblue!10,text=black] at (11.73,1.84) {GIVE};
\foreach \x in {11.05,12.42,13.45}{\sgnoexitsign[.92]{(\x,.86)}}  % 세 표지가 1.03 간격이라 원래 크기
\end{scope}

%% == 5. 벌금: 한꺼번에 뗌, B는 0으로 꺼짐 ==
\begin{scope}[xshift=\sgdE cm]
\node[fine] at (16.35,3.64) {$-$6,000};
\foreach \x in {15.4,16.35,17.3}
  \draw[sgastra,line width=.8pt,-{Stealth[length=4pt,width=3.5pt]}] (16.35,3.48)--(\x,2.86);
\sgrobot[\sgS]{(15.4,\sgY)}{sgblue}{}
\sgrobot[\sgS]{(16.35,\sgY)}{sggreen}{}
\sgrobotoff[\sgS]{(17.3,\sgY)}{}
\sgbar[\sgB]{(15.4,1.3)}{sgblue}{1}{4}   \sgcut[\sgB]{(15.4,1.3)}{2}{4}
\sgbar[\sgB]{(16.35,1.3)}{sggreen}{1}{4} \sgcut[\sgB]{(16.35,1.3)}{2}{4}
\sgbar[\sgB]{(17.3,1.3)}{sgoff}{0}{4}    \sgcut[\sgB]{(17.3,1.3)}{1}{4}
\node[font=\sffamily\bfseries\fontsize{8}{8}\selectfont,text=black,inner sep=0pt] at (17.3,1.02) {0};
\end{scope}

%% == 다음 라운드로 ==
\draw[sgblue!55,line width=.7pt,-{Stealth[length=4.5pt]}]
  (\sgW+.15,2.0)--(\sgW+.15,.3)--(-.15,.3)--(-.15,2.0)--(-.01,2.0);
\node[fill=white,sgfont=7,inner sep=2pt,text=black] at (\sgW/2,.3) {다음 라운드};
\node[tag,draw=sggold!80!black,fill=sggold!18,text=black,
  font=\sffamily\fontsize{6.8}{7}\selectfont] at (15.4+\sgdE,.3) {끝까지 남으면 상금 2,000만 $\div$ 남은 수};
\end{tikzpicture}%
